\chapter{How to search and screen}
\label{ch:procedure}

This is the procedure in the order it is carried out.
Each step names the command that does it and the chapter or section
that gives its grounds: a measurement of our own (\cref{ch:measured}),
the literature (\cref{ch:stopping,ch:abbreviated,ch:databases,%
ch:llm-screening,ch:reporting}), or the code.
The commands are those of this version of Scholar; \texttt{CLAIM} stands
for a session's name and \texttt{\dots} for what the reader fills in.

\begingroup\raggedright
\begin{description}
\item[Check the providers]
  \texttt{scholar providers check} probes each configured database with
  a one-result request and reports it as ok, rate limited, key rejected
  and so on (\cref{sec:provider-health}).
  A database that fails is not searched, and the method says so; a
  search that fails later is recorded as failed, never as zero hits.
  In health research no single database held what reviews included, and
  outside health the right databases differ (\cref{ch:databases}).

\item[One session per claim]
  Every query for one claim goes into one session,
  \texttt{scholar search \dots\ -n CLAIM}, so a work that several queries
  return is screened once.
  In a book chapter screened with one session per query, over a third of
  the screening was of works already screened in another session of the
  same claim (\cref{sec:measured-chapter}).

\item[Several wordings, fixed beforehand]
  Each concept of the claim gets several wordings (synonyms, the noun and
  the verb, the bare word beside the phrase), and each claim gets counter
  queries that reuse every concept term of the support queries.
  The wordings and the studies the queries must return
  (below) are written down before the first search.
  In the literature search of \cref{ch:lit-search}, the check by meaning
  found works the first wordings had missed, and new wordings of the
  same concepts brought several of them back.

\item[Every wording on every live provider, without a limit]
  \texttt{scholar search \textquotesingle{}\dots\textquotesingle{} -n CLAIM -{}-in title-abstract -p openalex
  -p s2 -p scopus -p wos -p ieee -p arxiv}, with no \texttt{-{}-limit}, so
  every result is fetched (\cref{sec:no-limit-default}).
  With \texttt{-{}-in title-abstract}, OpenAlex matches title and abstract
  instead of its default fields, which include the full text, and Scopus
  gets the query as \texttt{TITLE-ABS-KEY(\dots)}; Web of Science wraps a
  plain query as \texttt{TS=(\dots)}, its title, abstract and keyword
  fields (\cref{sec:providers-search-in}).
  A limit of 40 per query and database hid four of the sources a review
  later cited, one of them the only review against a hypothesis
  (\cref{sec:measured-limit}); what the literature says about shortened
  searches is in \cref{ch:abbreviated}, and the reporting guidelines ask
  that any limit be stated and justified (\cref{ch:reporting}).

\item[Check by meaning, then repair]
  \texttt{scholar sessions check CLAIM -s \textquotesingle{}<the claim as a question>\textquotesingle{}}
  looks each work of an OpenAlex semantic search up in the sessions
  (\cref{sec:sessions-check}).
  A work that bears on the claim and is in no session is a gap in the
  wordings: a new query by concept is added and run on every provider.
  In the literature search, 29 of the 50 works the check returned for one
  question were in no session, and two repair queries brought back 17 of
  them; of the works the checks returned that were in the sessions, 48 to
  73\,\% bore on the question, against 5 to 15\,\% of the first 100
  places of the keyword answers (\cref{sec:measured-tiers}).
  The check is a ranked list of at most 50 works, so it never counts
  anything and never shows that something does not exist.

\item[The named studies must come back]
  \texttt{scholar sessions check CLAIM -s \textquotesingle{}\dots\textquotesingle{} -k DOI} also locates
  each study named beforehand in the sessions.
  A query for a concept that returns none of the known studies on that
  concept is reworded and run again.

\item[Screen in tiers]
  The first tier is screened: the first 100 places of each query in each
  answer ordered by relevance (\texttt{sessions show} prints the order a
  search recorded), what the check by meaning returned, and the named
  studies.
  \texttt{scholar sessions export CLAIM -f screen-sheet -{}-pending
  -{}-top 100 -o BASE} writes the first part as a sheet and its counts to
  \texttt{BASE-tier.json} (\cref{sec:sheet-tier}); the verdicts go back
  with \texttt{scholar sessions decide CLAIM -{}-source agent
  -{}-reasons-file F} or are made in the review screens.
  A tier decides what is read first, not what exists: stopping a ranked
  screen at a fixed place or by a simple rule often misses relevant
  records (\cref{ch:stopping}), and relevant records sit deep in ranked
  lists (\cref{ch:databases}).

\item[Sample the rest]
  \texttt{scholar sessions export CLAIM -f screen-sheet -{}-pending
  -{}-top 100 -{}-rest-sample 60 -{}-seed S -o BASE} draws a seeded random
  sample of the pending records the tier leaves out and records the draw
  in the session (\cref{sec:rest-sample-sheet}); once it is screened,
  \texttt{scholar sessions tier-sample CLAIM -{}-report} gives the share of
  bearing records with its 95\,\% modified Wilson interval
  (\cref{sec:sessions-tier-sample,sec:llm-sample-interval}).
  That share is the one statement the method makes about what was not
  screened (\cref{sec:stopping-sample}).

\item[The classifier, and what its run line says]
  \texttt{scholar llm context CLAIM \textquotesingle{}\dots\textquotesingle{}} states the question and the
  criteria; \texttt{scholar llm classify CLAIM -{}-top 100 -n N} classifies
  the first tier's pending papers (\cref{sec:llm-classify-tier}).
  The run ends with one line of the form \enquote{Requests: 1 batch + 1
  retry of 2 papers; left pending: 0; tokens: \dots}: how many
  papers had to be asked again, how many were left without a usable
  answer, and the tokens the model reported (\cref{sec:classifier-run-cost}).
  In our test on 50 papers, runs with batches of different sizes agreed
  on 45 to 48 of the 50 kept-or-discarded verdicts, two runs with the
  same batch size on 49 (\cref{sec:measured-classifier}); what is known
  about screening with a
  language model is in \cref{ch:llm-screening}.

\item[A second screen where the result matters]
  \texttt{scholar llm classify CLAIM -{}-second-screen -{}-over machine}, or
  another screener's verdicts with \texttt{scholar sessions screen},
  stores a second verdict beside each decision without changing it
  (\cref{sec:llm-classify-second-screen,sec:sessions-screen});
  \texttt{scholar sessions agreement CLAIM} and \texttt{scholar sessions
  conflicts CLAIM -{}-unresolved} give the agreement and the papers to
  settle (\cref{sec:sessions-agreement,sec:sessions-conflicts}).
  One screener misses more relevant studies than two
  (\cref{sec:abbreviated-single}), and two runs of one model on the same
  papers did not give the same verdicts (\cref{ch:llm-screening}).

\item[Write what was done and what was not]
  The method states what was searched (databases, queries, dates, hits
  per query and database), what was screened (the tier's definition and
  size), what was not (how many records, and the sample's size, seed,
  count, rate and interval), and who screened.
  \texttt{scholar sessions show CLAIM} gives the hits per query and
  database, and \texttt{scholar sessions export CLAIM -f table
  -{}-bearing-only -{}-cited-from refs.bib} the table of records that bear on
  the claim (\cref{sec:audit-table,sec:audit-cited-from}).
  A negative result is worded to match: \enquote{none among the $N$
  screened of $M$ returned; a sample of $n$ of the rest held none
  (interval)}, never \enquote{none found}.
  The reporting guidelines behind this are in \cref{ch:reporting}.
\end{description}
\endgroup
